%==============================================================================
%  ClimaAtmos.jl — 1-hour hands-on tutorial
%  Companion to talks/main.tex (overview). This session is running & exploring
%  variables, not choosing physically relevant grids or timesteps.
%
%  COMPILE:  latexmk -pdf tutorial_1h.tex     (or pdflatex twice)
%  Participant scripts: tutorial_1h.jl, tutorial_1h_plots.jl
%==============================================================================
\documentclass[aspectratio=169,11pt]{beamer}

%------------------------------------------------------------------------------
% >>> THEME BLOCK <<<  (same as talks/main.tex)
%------------------------------------------------------------------------------
\usetheme{Madrid}
\usecolortheme{seahorse}
\definecolor{climablue}{RGB}{20,70,120}
\definecolor{climateal}{RGB}{20,130,130}
\setbeamercolor{structure}{fg=climablue}
\setbeamercolor{frametitle}{fg=white,bg=climablue}
\setbeamercolor{title}{fg=climablue}
\setbeamertemplate{navigation symbols}{}
\setbeamertemplate{caption}[numbered]
%------------------------------------------------------------------------------

\usepackage[utf8]{inputenc}
\usepackage[T1]{fontenc}
\usepackage{lmodern}
\usepackage{booktabs}
\usepackage{tikz}
\usetikzlibrary{positioning, arrows.meta}

\newcommand{\pkg}[1]{\texttt{#1}}
\newcommand{\cta}[1]{{\color{climateal}\textbf{#1}}}

\title[ClimaAtmos tutorial]{Running \texttt{ClimaAtmos.jl}}
\subtitle{A 1-hour hands-on: clone, instantiate, first runs, plots}
\author{Akshay Sridhar}
\institute{Climate Modeling Alliance (CliMA)}
\date{1-hour workshop}

\begin{document}

%==============================================================================
\begin{frame}[plain]
  \titlepage
  \begin{center}\small
    Goal: a working clone, two short simulations, and a look at the output.\\
    Not today: production resolution, CFL, or physically tuned cases.
  \end{center}
\end{frame}

%==============================================================================
\begin{frame}{This hour}
\begin{center}
\begin{tabular}{clp{0.55\textwidth}}
  \toprule
  \textbf{min} & \textbf{block} & \textbf{you will} \\
  \midrule
  0--5   & Welcome        & know what we will / will not do \\
  5--15  & Clone          & have the repository on disk \\
  15--25 & Instantiate    & have a compiled Julia environment \\
  25--45 & First runs     & run a column and a sphere \\
  45--55 & Visualize      & plot diagnostics from NetCDF \\
  55--60 & Recap          & leave with a cheat sheet \\
  \bottomrule
\end{tabular}
\end{center}
\vspace{0.6em}
\small
Keep \textbf{one} Julia session open. The first compile is slow; restarting
throws that work away.
\note{If instantiate was done as homework, spend the recovered 10 min on
inspecting \texttt{Y} and extra plots. If it was not, start instantiate
immediately after clone and talk through the mental model while it runs.}
\end{frame}

%==============================================================================
\begin{frame}{What this is not}
\begin{itemize}
  \item Not a dynamics lecture --- the overview talk covers that.
  \item Not a resolution / timestep / physics-tuning session.
    The grids here are \cta{coarse on purpose} so they finish in the room.
  \item Not ``BOMEX as published'' or ``baroclinic wave at day 10''.
    We will see the \emph{variables} those cases write, on a laptop budget.
\end{itemize}
\vspace{0.8em}
Three objects, used all afternoon:
\begin{itemize}
  \item \texttt{Y} --- prognostic state (\texttt{Y.c} centres, \texttt{Y.f} faces)
  \item \texttt{p} --- cache of precomputed fields and parameters
  \item diagnostics --- derived NetCDF output (\texttt{ta}, \texttt{ua}, \texttt{hus}, \ldots)
\end{itemize}
\note{Write \texttt{Y}/\texttt{p} on the board. Centre vs face is the only
discretization fact they need today.}
\end{frame}

%==============================================================================
\begin{frame}[fragile]{0. Prerequisites}
Before (or while) we start:
\begin{itemize}
  \item \textbf{Julia 1.11} via \href{https://julialang.org/install/}{juliaup}
  \item \textbf{Git}, and \textbf{$\sim$10\,GB} disk (deps + artifacts)
  \item Network for the first instantiate
\end{itemize}
\vspace{0.4em}
\begin{semiverbatim}
curl -fsSL https://install.julialang.org | sh
juliaup add 1.11
juliaup default 1.11
julia --version          # 1.11.x
\end{semiverbatim}
\vspace{0.4em}
\small
Optional but useful, in the \emph{default} environment (plain \texttt{julia},
no \texttt{--project}):
\begin{semiverbatim}
julia -e 'using Pkg; Pkg.add(["ClimaAnalysis","CairoMakie"])'
\end{semiverbatim}
\note{Do not Pkg.add plotting packages into the ClimaAtmos project --- that
edits Project.toml. Same rule as CUDA in the installation docs.}
\end{frame}

%==============================================================================
\begin{frame}[fragile]{1. Clone}
HTTPS is the reliable workshop default (no SSH keys required):
\vspace{0.4em}
\begin{semiverbatim}
git clone https://github.com/CliMA/ClimaAtmos.jl.git
cd ClimaAtmos.jl
ls
\end{semiverbatim}
\vspace{0.5em}
What you should see, at a glance:
\begin{itemize}
  \item \texttt{src/} --- the model
  \item \texttt{config/} --- YAML cases (\texttt{model\_configs/}, \texttt{default\_configs/})
  \item \texttt{talks/} --- this tutorial's scripts
  \item \texttt{Project.toml} / \texttt{Manifest.toml} --- pinned environment
\end{itemize}
\small
Docs: \url{https://clima.github.io/ClimaAtmos.jl/dev/installation/}
\note{Walk the tree for 60 seconds. Point at config/model\_configs so YAML
later is not a surprise. Do not dive into src/.}
\end{frame}

%==============================================================================
\begin{frame}[fragile]{2. Instantiate (precompile)}
From the repository root:
\vspace{0.3em}
\begin{semiverbatim}
julia +1.11 --project
\end{semiverbatim}
Then in the REPL:
\begin{semiverbatim}
using Pkg
Pkg.instantiate()    # download + precompile the Manifest
import ClimaAtmos as CA
\end{semiverbatim}
\vspace{0.4em}
\begin{block}{Expect a wait}
The first instantiate downloads and precompiles the full stack and can take
\textbf{tens of minutes}. Later sessions start quickly. Leave this REPL open.
\end{block}
\small
Equivalent one-liner: \texttt{julia +1.11 --project -e 'using Pkg; Pkg.instantiate()'}
\note{Start this before explaining Y if the room is cold. Filler: two
interfaces (script vs YAML), output layout, short-name diagnostics.}
\end{frame}

%==============================================================================
\begin{frame}{While that compiles}
Two equivalent ways to launch a run --- we will use the \cta{script} path:
\vspace{0.4em}
\begin{center}
\begin{tabular}{lll}
  \toprule
   & \textbf{script} & \textbf{YAML} \\
  \midrule
  entry & \texttt{AtmosSimulation} / \texttt{Presets} & \texttt{AtmosConfig("file.yml")} \\
  best for & REPL, notebooks, this hour & CI, sharing a case \\
  \bottomrule
\end{tabular}
\end{center}
\vspace{0.8em}
Output always lands under \texttt{output/<job\_id>/}:
\begin{itemize}
  \item numbered run directories \texttt{output\_0000}, \texttt{output\_0001}, \ldots
  \item \texttt{output\_active} $\to$ the latest
  \item \textbf{NetCDF} (\texttt{.nc}) = diagnostics; \textbf{HDF5} = checkpoints
\end{itemize}
\end{frame}

%==============================================================================
\begin{frame}[fragile]{Diagnostics we can plot}
Default diagnostics are time averages, and they write nothing for runs
shorter than 1 hour (at 1 hour you get a single hourly mean). Request
instantaneous snapshots so we can see evolution --- and skip \texttt{hus}
on dry cases:
\vspace{0.3em}
\begin{semiverbatim}
diag = CA.DiagnosticsConfig(;
    default = false,
    additional = (
        "ta"   => (; period = "10mins"),
        "hus"  => (; period = "10mins"),
        "ua"   => (; period = "10mins"),
        "rhoa" => (; period = "10mins"),
    ),
)
\end{semiverbatim}
\small
Diagnostics are fixed at construction. For the dry wave, drop \texttt{hus}
and use a 6-hour period (\texttt{sphere\_diagnostics} in the script).
\end{frame}

%==============================================================================
\begin{frame}[fragile]{3a. Single column --- BOMEX}
Shallow-cumulus \emph{initial state} in one column. The preset is moist 0-moment
microphysics, \textbf{no} EDMF: enough physics to have humidity, cheap enough
to finish here.
\vspace{0.3em}
\begin{semiverbatim}
import ClimaAtmos as CA

column = CA.Presets.bomex(
    Float32;
    t_end = "1hours",
    job_id = "workshop_bomex",
    diagnostics = diag,
)
\end{semiverbatim}
\vspace{0.3em}
\small
Defaults inside the preset: 60 uniform levels, $z_\mathrm{max} = 3$\,km,
\texttt{dt = 10}\,s. Construction fills \texttt{Y} but does \emph{not} step time.
\note{Say out loud: this is not the published BOMEX + PROPHET YAML. That
config is prognostic\_edmfx\_bomex\_column.yml; we will name it, not run it.}
\end{frame}

%==============================================================================
\begin{frame}[fragile]{Inspect \texttt{Y} before running}
\begin{semiverbatim}
Y = column.integrator.u
propertynames(Y.c)   # centres: ρ, uₕ, ρe_tot, ρq_tot, ...
propertynames(Y.f)   # faces:   u₃, ...
extrema(Y.c.ρ)
extrema(Y.c.ρe_tot)
\end{semiverbatim}
\vspace{0.5em}
\begin{itemize}
  \item Prognostic variables are densities ($\rho$, $\rho e_\mathrm{tot}$,
        $\rho q_\mathrm{tot}$) plus velocity.
  \item Diagnostics later are the familiar names: \texttt{ta}, \texttt{hus},
        \texttt{ua}, \ldots
  \item \texttt{p = column.integrator.p} is the cache --- peek, don't mutate.
\end{itemize}
\note{Have people shout out propertynames. If compile is still going, this
slide is the waiting-room activity once AtmosSimulation returns.}
\end{frame}

%==============================================================================
\begin{frame}[fragile]{Run the column}
\begin{semiverbatim}
CA.solve_atmos!(column)
column.integrator.t      # should be 3600 s
column.output_dir        # output/workshop_bomex/output_0000
\end{semiverbatim}
\vspace{0.5em}
First \texttt{solve\_atmos!} in a session compiles the time stepper and can
take several minutes. The 1-hour integration itself is short on a column.
\vspace{0.5em}
Optional: one timestep at a time
\begin{semiverbatim}
import ClimaTimeSteppers: step!
step!(column.integrator)
\end{semiverbatim}
\end{frame}

%==============================================================================
\begin{frame}[fragile]{3b. Sphere --- dry baroclinic wave}
Idealized mid-latitude jet + a perturbation, on a cubed sphere. Dry, no
surface fluxes. Default grid: \texttt{h\_elem = 6} ($\sim$550\,km),
\texttt{z\_elem = 10}, \texttt{dt = 10}\,min.
\vspace{0.3em}
\begin{semiverbatim}
sphere = CA.Presets.baroclinic_wave(
    Float32;
    t_end = "1days",
    job_id = "workshop_bcw",
    diagnostics = diag_dry,   # like diag, but no hus; 6-hour period
)
Y = sphere.integrator.u
propertynames(Y.c)     # no ρq_tot
CA.solve_atmos!(sphere)
\end{semiverbatim}
\small
A published wave is typically integrated 8--10 days. One day is enough to
write fields and look at them. Do not read this as a converged test.
\note{Show that moist vs dry is visible in propertynames. hus on a dry model
errors --- that is a useful live demo if someone tries it.}
\end{frame}

%==============================================================================
\begin{frame}{What you just ran}
\begin{center}
\begin{tabular}{lll}
  \toprule
   & \textbf{column} & \textbf{sphere} \\
  \midrule
  preset & \texttt{Presets.bomex} & \texttt{Presets.baroclinic\_wave} \\
  grid & 60 levels, 3\,km & $h=6$, 10 levels, 30\,km \\
  physics & moist 0M, no EDMF & dry, no surface flux \\
  model time & 1 hour & 1 day \\
  look at & \texttt{ta}, \texttt{hus} vs $z$ & \texttt{ua} zonal mean \\
  \bottomrule
\end{tabular}
\end{center}
\vspace{0.8em}
Both used \texttt{Float32}. Output is under
\texttt{output/workshop\_bomex/} and \texttt{output/workshop\_bcw/}.
\end{frame}

%==============================================================================
\begin{frame}[fragile]{4. Visualize --- a second Julia}
Plotting packages go in the \cta{default} environment, not the model project.
\vspace{0.3em}
\begin{semiverbatim}
# new terminal, from ClimaAtmos.jl, *no* --project
julia +1.11
\end{semiverbatim}
\begin{semiverbatim}
import Pkg; Pkg.add(["ClimaAnalysis", "CairoMakie"])
using ClimaAnalysis
import ClimaAnalysis.Visualize as viz
import CairoMakie

simdir = SimDir("output/workshop_bomex/output_active")
ta = get(simdir, "ta")
\end{semiverbatim}
\small
\texttt{SimDir} catalogs every diagnostic the run wrote. \texttt{get} loads
one by short name as an \texttt{OutputVar}.
\end{frame}

%==============================================================================
\begin{frame}[fragile]{Column profiles}
\begin{semiverbatim}
using ClimaAnalysis: times
t0 = first(times(ta))
hus = get(simdir, "hus")

fig = CairoMakie.Figure(; size = (900, 400))
viz.plot!(fig[1, 1], ta;  time = t0, lon = 0, lat = 0)
viz.plot!(fig[1, 2], hus; time = t0, lon = 0, lat = 0)
CairoMakie.save("workshop_bomex_ic.png", fig)
\end{semiverbatim}
\vspace{0.4em}
\texttt{viz.plot!} picks a line or a heatmap from how many dimensions remain.
Drop extra axes with \texttt{slice} / keyword slices; reduce with
\texttt{average\_lon}, \texttt{average\_time}, \ldots
\end{frame}

%==============================================================================
\begin{frame}[fragile]{Sphere fields}
\begin{semiverbatim}
bcw = SimDir("output/workshop_bcw/output_active")
ua = get(bcw, "ua")
ta = get(bcw, "ta")
t = last(times(ua))

fig = CairoMakie.Figure()
viz.plot!(fig, average_lon(ua); time = t)   # lat--height
CairoMakie.save("workshop_bcw_ua_zonal_mean.png", fig)

fig = CairoMakie.Figure()
viz.plot!(fig, ta; z = 0, time = t)         # lat--lon near surface
CairoMakie.save("workshop_bcw_ta_surface.png", fig)
\end{semiverbatim}
\small
Full copy-paste: \texttt{talks/tutorial\_1h\_plots.jl}
\end{frame}

%==============================================================================
\begin{frame}[fragile]{YAML path (optional)}
The same runs, file-driven --- later files override earlier keys:
\vspace{0.3em}
\begin{semiverbatim}
config = CA.AtmosConfig(
    ["config/model_configs/baroclinic_wave.yml",
     "talks/workshop_configs/short_sphere.yml"];
    job_id = "workshop_bcw_yaml",
)
sim = CA.AtmosSimulation(config)
CA.solve_atmos!(sim)
\end{semiverbatim}
\vspace{0.3em}
Or edit one key in the REPL after loading:
\begin{semiverbatim}
config = CA.AtmosConfig("config/model_configs/baroclinic_wave.yml")
config.parsed_args["t_end"] = "1days"
\end{semiverbatim}
\small
The published BOMEX+PROPHET case is
\texttt{config/model\_configs/prognostic\_edmfx\_bomex\_column.yml} ---
that is the physically complete column, not what we ran live.
\end{frame}

%==============================================================================
\begin{frame}{Cheat sheet}
{\small
\begin{tabular}{@{}ll@{}}
  \toprule
  clone & \texttt{git clone https://github.com/CliMA/ClimaAtmos.jl.git} \\
  env & \texttt{julia +1.11 --project} then \texttt{Pkg.instantiate()} \\
  column & \texttt{CA.Presets.bomex(Float32; t\_end="1hours", job\_id=...)} \\
  sphere & \texttt{CA.Presets.baroclinic\_wave(Float32; t\_end="1days", ...)} \\
  run & \texttt{CA.solve\_atmos!(simulation)} \\
  state & \texttt{Y = simulation.integrator.u} \\
  output & \texttt{output/<job\_id>/output\_active} \\
  plot env & default Julia + \texttt{ClimaAnalysis}, \texttt{CairoMakie} \\
  load & \texttt{SimDir(...) ; get(simdir, "ta")} \\
  \bottomrule
\end{tabular}
}
\vspace{0.6em}
Scripts in this repo: \texttt{talks/tutorial\_1h.jl}, \texttt{talks/tutorial\_1h\_plots.jl}
\end{frame}

%==============================================================================
\begin{frame}{Where to go next}
\begin{itemize}
  \item Docs: \url{https://clima.github.io/ClimaAtmos.jl/dev/}
        --- first simulation, single-column cases, visualizing output
  \item Next workshop hour (not today): grids at physically relevant scales,
        timesteps / CFL, PROPHET, radiation
  \item YAML library: \texttt{config/model\_configs/} and
        \texttt{config/longrun\_configs/}
\end{itemize}
\vspace{0.8em}
\begin{center}
\large Questions --- while we still have a live REPL.
\end{center}
\end{frame}

%==============================================================================
\appendix
\begin{frame}[fragile]{Backup: if instantiate fails}
Typical recovery, still in \texttt{julia --project}:
\begin{semiverbatim}
using Pkg
Pkg.resolve()
Pkg.instantiate()
\end{semiverbatim}
\vspace{0.4em}
If a download stalled, retry on a better network. Do not \texttt{Pkg.update()}
in a workshop --- that floats off the Manifest.
\vspace{0.4em}
Shared depot (instructors): set \texttt{JULIA\_DEPOT\_PATH} to a pre-warmed
\texttt{.julia} so the room skips the download.
\end{frame}

\begin{frame}[fragile]{Backup: inspecting a \texttt{Field} without plots}
\begin{semiverbatim}
import ClimaCore as CC
z = CC.Fields.coordinate_field(Y.c).z
extrema(Y.c.ρ)
any(isnan, parent(Y.c.ρ))
parent(Y.c.ρ)          # raw array, column is small
\end{semiverbatim}
Useful when CairoMakie is not installed yet, or when a run looks wrong.
\end{frame}

\begin{frame}{Backup: diagnostic short names we used}
\begin{center}
\begin{tabular}{lll}
  \toprule
  name & quantity & notes \\
  \midrule
  \texttt{ta} & air temperature & both cases \\
  \texttt{thetaa} & potential temperature & column \\
  \texttt{hus} & specific humidity & moist only \\
  \texttt{ua},\texttt{va},\texttt{wa} & wind components & \\
  \texttt{rhoa} & air density & \\
  \texttt{pfull} & pressure & sphere \\
  \bottomrule
\end{tabular}
\end{center}
\vspace{0.5em}
Catalog: docs $\to$ \emph{Available diagnostic variables}. Requesting
\texttt{hus} on a dry model errors at setup --- that is expected.
\end{frame}

\end{document}
